% Report section: "Results and cost" (~three quarters of a page).
% Self-contained fragment; \input this into the main report. Needs booktabs.
% Every number is from outputs/results/{results.csv,calls.csv} and outputs/traces/ (run 20261013-200839).

\section{Results and cost}

All 90 task-runs completed: three full repeats of T1--T10 for each condition.
Table~\ref{tab:first-try} shows first-attempt success per task.

\begin{table}[h]
\centering
\small
\begin{tabular}{llccc}
\toprule
Task & Lessons needed & No memory & Episodic & Structured \\
\midrule
T1  & C1      & 0/3 & 0/3 & 1/3 \\
T2  & C1, C2  & 0/3 & 0/3 & 0/3 \\
T3  & C1, C2  & 0/3 & 3/3 & 0/3 \\
T4  & C1, C2  & 0/3 & 2/3 & 3/3 \\
T5  & C1, E1  & 2/3 & 0/3 & 2/3 \\
T6  & C1, E2  & 3/3 & 3/3 & 3/3 \\
T7  & C1, C2  & 0/3 & 3/3 & 3/3 \\
T8  & C2, E1  & 0/3 & 3/3 & 1/3 \\
T9  & E1, E2  & 3/3 & 3/3 & 3/3 \\
T10 & C1, C2  & 0/3 & 3/3 & 3/3 \\
\midrule
First-attempt success, mean (sd) over repeats & & 26.7\% (5.8) & 66.7\% (5.8) & 63.3\% (5.8) \\
Final success within 3 attempts & & 28/30 & 30/30 & 30/30 \\
Attempts per task, mean & & 2.13 & 1.33 & 1.40 \\
\bottomrule
\end{tabular}
\caption{First-attempt successes out of 3 repeats per task and condition.}
\label{tab:first-try}
\end{table}

\paragraph{Hypotheses.}
\textbf{H1} (on T2--T4, T7, T10): no memory stayed at or below 30\% (0/15) --- \emph{held}.
Structured memory reached 9/15 (60.0\%), below the 80\% threshold --- \emph{failed}; it did
use fewer attempts than no memory on these tasks (1.47 vs.\ 2.67) --- \emph{held}. Episodic
memory, which H1 makes no claim about, reached 11/15 (73.3\%).
\textbf{H2}: on T5, episodic memory was worse than no memory (0/3 vs.\ 2/3), and all three
episodic failures were the planted 08:00-shift-on-Wingra mistake --- \emph{held} for episodic;
structured memory tied with no memory (2/3) --- \emph{failed} for structured. On T6, episodic
and structured tied at 3/3 --- \emph{failed}. On T1, where memory is empty in every condition,
results were 0/3, 0/3 and 1/3 --- \emph{held} up to one run of variation.

\paragraph{Surprising comparisons.}
Episodic memory beat structured memory on the rule-reuse tasks, mainly on T3 (3/3 vs.\ 0/3)
and T8 (3/3 vs.\ 1/3). On T3, structured memory retrieved the 08:00 day rule in all three
repeats, yet the agent used the calendar day every time; the memory logged this as
``retrieved but not applied''. On T8, structured memory succeeded first try only in the one
repeat where it had already stored the Wingra exception at T5. Overall first-attempt success
was stable across repeats (no memory 20/30/30\%, episodic 70/70/60\%, structured
60/70/60\%), but each per-task cell rests on only three runs, and we did not run extra
repeats to check these comparisons.

\paragraph{Cost.}
The experiment made 452 model calls (583{,}723 input and 35{,}338 output tokens) costing
\$0.159, or 3.2\% of the \$5 limit: no memory 202 calls (\$0.080), episodic 105 (\$0.030),
structured 145 (\$0.049). All were agent calls; neither memory system calls the model.
Including our development test runs, \texttt{calls.csv} records 555 calls costing \$0.193.

\paragraph{Limitation.}
Six wrong answers, all in the no-memory condition (on T5, T7, T8 twice and T10 twice),
matched no planted mistake and received the checker's generic message with no hint; both
no-memory T10 runs that received one failed all three attempts, so the no-memory baseline is
likely understated on T10.
